\documentclass{article}
\usepackage[utf8]{inputenc}
\usepackage{amsmath}
\usepackage{amssymb}

% KantiKoala LaTeX Unterstützung
% ==============================
% Laden Sie dieses Dokument hoch, um zu sehen, wie es konvertiert wird.

\begin{document}

\section{Einführung in LaTeX}

Dieses Dokument zeigt, welche LaTeX-Befehle von KantiKoala direkt unterstützt werden.
Der Text wird automatisch formatiert.

\subsection{Unterstützte Formatierungen}

Sie können \textbf{fetten Text}, \textit{kursiven Text} oder \underline{unterstrichenen Text} verwenden.

\subsection{Listen}

Normale Aufzählungen funktionieren:

\begin{itemize}
    \item Erster Punkt
    \item Zweiter Punkt mit \textbf{wichtigem} Detail
\end{itemize}

Nummerierte Listen funktionieren ebenfalls:

\begin{enumerate}
    \item Schritt Eins
    \item Schritt Zwei
\end{enumerate}

\section{Mathematik}

KantiKoala nutzt \textbf{KaTeX} für die Darstellung von Formeln. Alle gängigen Mathe-Pakete (amsmath, amssymb) werden unterstützt.

Inline-Formeln: $f(x) = x^2$ oder $e^{i\pi} + 1 = 0$.

Abgesetzte Formeln:

$$ \int_{-\infty}^{\infty} e^{-x^2} \, dx = \sqrt{\pi} $$

Oder komplexere Brüche:
$$ x = \frac{-b \pm \sqrt{b^2 - 4ac}}{2a} $$

\section{Hinweise für Lehrpersonen}

\begin{itemize}
    \item Verwenden Sie \texttt{section} und \texttt{subsection} für die Gliederung.
    \item Bilder werden aktuell in LaTeX-Dateien NICHT unterstützt. Bitte nutzen Sie dafür Word (.docx).
    \item Komplizierte Tabellen sollten vermieden werden oder als Bild eingefügt werden.
    \item TikZ-Grafiken werden nicht gerendert.
\end{itemize}

\end{document}
